\section{Introduction}
\label{sec:introduction}

Electromagnetic spectrum monitoring systems are expected to recognize known activity while operating across changing emitters, receivers, sites, propagation conditions, acquisition schedules, and interference regimes. Conventional supervised recognition assumes that evaluation samples belong to the same label and domain support as the training set. That assumption is fragile in open-spectrum operation, where an unknown signal technology, a new acquisition domain, or a previously unseen jammer scenario may be encountered without an OOD label available at decision time. Closed-set models can remain confident under such shifts, so classification confidence and OOD separability must be evaluated as related but distinct properties. \pendingcite{R1} \pendingcite{R2}

RF OOD evaluation also has a protocol-design problem. An apparent detector can be improved after the fact by reversing a score, selecting the strongest test comparator, fitting normalization on evaluation samples, or optimizing a threshold using OOD labels. Those operations leak information from the target shift and can turn a diagnostic observation into an overstated primary result. A deployment-relevant experiment instead fixes the score direction and analysis roles in advance, fits classifiers and distance models on training data, fits calibration and normalization on ID validation data, and leaves the test ID and test OOD rows untouched until evaluation. \pendingcite{R3}

This work studies that disciplined path on the OpenEW-SA artifact convention. The term \emph{multi-view} has two scopes here. Across datasets, the evaluation includes processed power-spectral-density (PSD), in-phase/quadrature (I/Q), and tabular radio/network representations. Within the proposed detector, the fused views are complementary evidence sources: predictive uncertainty and two forms of feature-space distance. The current experiments do not claim an end-to-end, within-sample fusion of raw I/Q, PSD, spectrogram, and metadata inputs; that broader multimodal model remains future work.

The study makes four contributions. First, it defines frozen class-OOD and domain/scenario-OOD evaluations for ElectroSense, DeepSense, and JamShield while preserving symbolic RF labels. Second, it connects a lightweight supervised recognition baseline, validation-fitted temperature scaling, and train-fitted feature-distance scores in one reproducible pipeline. Third, it specifies \method{ts_entropy_cosine_euclidean} as an equal-weight primary fusion after ID-validation-only robust normalization, while keeping the Mahalanobis-augmented fusion exploratory. Fourth, it reports fixed-orientation point estimates and paired bootstrap uncertainty, including a negative DeepSense result and a post-hoc inversion diagnostic that is explicitly excluded from the primary analysis.

The resulting evidence is intentionally bounded. The primary method performs differently across the three shifts, and no universal generalization or universal statistical-significance claim is made. Instead, the analysis asks where calibrated uncertainty and feature geometry agree, where they trade off across metrics, and where their assumed higher-is-more-OOD direction breaks under domain change.
